\section{Experimental Setup}

This section specifies the planned experimental protocol. Configuration choices marked as proposed will be resolved using training and validation data before held-out evaluation. No retrieval results are implied by the following setup.

\subsection{Dataset}

\textbf{Budgeted retrieval pilot.} We select MSR-VTT \cite{xu2016msrvtt} as the source and plan a custom sample of 500 unique videos, rather than processing the full dataset. This sample is intended to test the fusion hypothesis within a limited computational budget. It is not the standard 1K-A protocol, and its scores will not be compared numerically with full-benchmark results. MSR-VTT provides reference descriptions supporting query-to-video relevance; actual video availability and annotation integrity will be verified before sampling.

\textbf{Sampling and eligibility.} Before any retrieval scores are inspected, we build an eligible manifest of available, decodable videos with valid IDs and at least one nonempty reference description. We report the source release, eligible population, exclusions, and download failures. From the sorted manifest, videos are sampled without replacement using seed 42. When source category labels are available, proportional stratification uses largest-remainder quotas; otherwise uniform random sampling is used and the realized category distribution is reported. Semantic query categories used for fusion are distinct from source dataset categories. Clips are not selected according to branch performance, caption quality, or apparent retrieval difficulty.

\textbf{Video-disjoint supervision.} The sampled video IDs are assigned to 300 training, 100 validation, and 100 test groups, preserving source-category proportions when feasible. All descriptions of a video remain in the same supervision partition. Within each video, at most five distinct reference descriptions are sampled using a fixed seed before model fitting. This yields at most 1,500 training, 500 validation, and 500 test queries; actual counts are reported. The descriptions are query annotations only and are excluded from the generated-caption index. Duplicate text associated with different videos is flagged for relevance review and excluded from query sampling if ambiguous; it is not automatically labeled as a negative for another matching video.

\textbf{Shared candidate collection.} Every query ranks the same 500 candidate videos. Training, validation, and test refer to query-supervision partitions; they do not partition the retrieval index. Frozen visual and caption features for all candidates are computed independently of the query annotations. Test-positive videos can therefore appear as distractors for other query groups, but test queries and their relevance labels never enter gate training or model selection. This is a fixed-collection pilot, not a claim of generalization to an unseen candidate collection. Learned video projections or encoders are outside this protocol and would require a separate candidate-disjoint design.

\textbf{Index construction.} We initially sample four frames per video at relative positions (i + 0.5)/4 for i = 0,1,2,3 and generate one caption per frame using a frozen model without query input. The default budget therefore contains 2,000 frame-caption pairs. A fixed decoding policy handles short or partially decodable clips, and duplicate frame timestamps are logged. Each branch takes the maximum query similarity across its four indexed items to produce one video score. All methods reuse the identical frame, caption, and feature caches.

\textbf{Runtime feasibility.} Before processing the full sample, a fixed 20-video subset drawn from the training partition measures loading, frame extraction, caption generation, and encoding time. Warm-up and model-download time are recorded separately. These videos retain their training assignment. The runtime estimate is T\_setup + N * t\_video + T\_fusion/evaluation, where t\_video is measured under the intended batch size. The estimate guides the budget but is not a promise of completion within one day. If the measured budget is unsuitable, a prespecified reduced tier uses 200 videos with 120/40/40 training/validation/test groups and the same four-frame rule. The tier is chosen before validation/test retrieval outcomes are inspected; the 200-video collection is a nested subset preserving the original group assignments. Results from different candidate-pool sizes are reported separately.

\textbf{Scope of evidence.} With 100 test-video groups in the default pilot, multiple descriptions increase query coverage but not the number of independent video groups. Confidence intervals therefore resample video groups. Conclusions are restricted to the sampled population and candidate size; full-dataset generalization, cross-domain robustness, and benchmark superiority remain untested. The previous broadcast collection may be evaluated separately once its data and relevance protocol are verified.

\subsection{Baselines}

All primary methods use identical frozen branch outputs and the same candidate collection. We plan the following comparisons:

\begin{center}
\small
\begin{tabularx}{\textwidth}{@{}lXXX@{}}
\toprule
ID & Method & Weighting rule & Purpose \\
\midrule
B1 & Visual-only & Visual branch alone & Establish visual retrieval performance \\
B2 & Caption-only & Generated-caption branch alone & Measure standalone caption retrieval \\
B3 & Equal-weight RRF & Equal contribution from both branches & Test untuned combination \\
B4 & Globally tuned RRF & One coefficient for every query & Strong fixed fusion control \\
B5 & Semantic-category RRF & Category-specific coefficients & Interpretable query adaptation \\
B6 & Linear query gate & Affine query projection followed by softmax & Isolate the benefit of query conditioning without a hidden layer \\
P & MLP query gate & Nonlinear query projection followed by softmax & Candidate learned method \\
\bottomrule
\end{tabularx}
\end{center}

For B3--P, the fused score is

\[
S(q,v)=w_I(q)b_I(q,v)+w_C(q)b_C(q,v),\qquad w_I(q)+w_C(q)=1,
\]

where \(b_m(q,v)=1/(k+r_m(q,v))\) if video \(v\) is present in branch \(m\)'s retrieved list and zero otherwise. Ranks are one-based. B4 uses a constant weight, while B5--P condition the weights on the query. Each method sorts the resulting scores using the same deterministic video-ID tie rule. Full rankings are used for the selected 500-video or 200-video pilot so that truncation does not confound fusion effectiveness.

B5 predicts object/attribute, action/interaction, or scene/other categories using a frozen query-only classifier. Category weights are estimated using development data, with shrinkage toward the global coefficient when support is limited. The classifier, label policy, and shrinkage strength will be fixed before test evaluation. Manual semantic labels, if collected, will be used for diagnostic analysis and will not silently replace the deployed classifier.

B6 and P receive the same normalized query embedding and the same training targets. The comparison therefore isolates the extra nonlinear capacity of the MLP. An oracle coefficient selected using each test query's relevance labels may be reported separately as a diagnostic ceiling over the declared weight grid. It is not a deployable method or a held-out prediction, and it cannot be used to select the final model.

\subsection{Metrics}

For a query \(q\), let \(R(q)\) be its relevant-video set and \(T_K(q)\) the top-K retrieved videos. We define

\[
\operatorname{Recall@K}=\frac{1}{|Q|}\sum_{q\in Q}
\frac{|T_K(q)\cap R(q)|}{|R(q)|},
\]

\[
\operatorname{Hit@K}=\frac{1}{|Q|}\sum_{q\in Q}
\mathbb{1}[T_K(q)\cap R(q)\ne\varnothing].
\]

These metrics coincide for single-positive retrieval. We report K \(\in\) \{1,5,10\}, naming the metric according to the verified relevance protocol. We additionally report mean reciprocal rank,

\[
\operatorname{MRR}=\frac{1}{|Q|}\sum_{q\in Q}\frac{1}{r_q^*},
\]

where \(r_q^*\) is the rank of the first relevant video and the contribution is zero if no relevant video is returned. Full-list MRR is the proposed primary selection metric. If MRR@10 is needed for a historical comparison, ranks beyond ten contribute zero and the truncated metric is labeled explicitly.

The primary comparison is P versus B4. We report paired per-query differences and 95\% bootstrap confidence intervals, using 10,000 resamples and grouping queries by video when they share relevance targets. For learned models, we additionally report mean and standard deviation across training seeds rather than treating a bootstrap interval from one run as training uncertainty. Comparisons with the other methods are exploratory unless separately prespecified; multiple formal significance claims will require correction.

To characterize behavior beyond aggregate metrics, we report semantic-category performance, per-query wins/ties/losses, and the distribution of predicted weights. Online efficiency is measured by median and 95th-percentile query latency, including query encoding, weight prediction, branch retrieval, and fusion. Offline caption generation and video encoding are reported separately. All latency comparisons use the same hardware, candidate pool, and warm-up protocol.

\subsection{Configurations}

\textbf{Frozen retrieval pipeline.} Visual and caption encoders will remain frozen for the main comparison. Similarity is computed only between embeddings produced by a verified compatible query--candidate encoder pair. Different branches need not share vector dimensions because their rankings are fused. Exact checkpoint IDs, revisions, preprocessing, tokenization, pooling, embedding dimensions, caption prompt, and decoding settings are pending inspection of the implementation. [TODO: populate from code and run metadata; do not infer checkpoints from model-family names.]

\textbf{Fusion settings.} Proposed defaults are k = 60 and full rankings for the selected pilot candidate pool. The fixed-weight search uses G = \{0,0.1,\ldots{},1\}, optimized on training/development data with validation-based selection as appropriate to the chosen protocol. A prespecified sensitivity analysis considers k \(\in\) \{10,30,60,100\}, with matching validation budgets for fixed and adaptive methods. Ties between equally effective coefficients favor the global coefficient when available, then proximity to 0.5, and finally the smaller coefficient. No tie rule uses test outcomes.

\textbf{MLP gate.} The initial architecture is Linear(d,32) \(\rightarrow\) GELU \(\rightarrow\) Dropout(0.1) \(\rightarrow\) Linear(32,2) \(\rightarrow\) Softmax, where d is the verified query embedding dimension. The two outputs are visual and caption weights. Only the gate is trainable. The linear baseline removes the hidden layer and uses the same query input and loss.

The proposed training target for a query is the coefficient in G that gives the best retrieval utility using that training query's relevance annotations. Ties are resolved toward a global coefficient estimated on training data, then by a deterministic rule. Queries with identical utility for every coefficient provide no branch-preference signal and are excluded from the target-regression loss; their number is reported. The initial loss is mean squared error between the predicted visual weight and the target coefficient. This objective avoids differentiating through the rank operation; alternative ranking losses will be evaluated separately rather than mixed into the initial gate comparison.

\begin{center}
\small
\begin{tabularx}{\textwidth}{@{}XX@{}}
\toprule
Setting & Proposed initial value or policy \\
\midrule
Query input & L2-normalized frozen text embedding; checkpoint pending \\
MLP hidden dimension & 32 \\
Dropout & 0.1 \\
Optimizer & AdamW \\
Learning rate & 1e-3 \\
Weight decay & 1e-4 \\
Batch size & 32 \\
Maximum epochs & 100 \\
Checkpoint selection & Highest validation MRR; earliest epoch on ties \\
Early stopping & 10 epochs without validation MRR improvement \\
Training seeds & 13, 42, 2026; report all, never select the best seed \\
Split seed & 42; video-group assignments shared across methods \\
Default sample & 500 videos; 300/100/100 supervision groups \\
Reduced tier & 200 videos; 120/40/40 supervision groups \\
Frames per video & 4; one independently generated caption per frame \\
Query cap & 5 distinct reference descriptions per video \\
RRF offset & 60 \\
Weight grid & \{0,0.1,\ldots{},1\} \\
Hardware and software & TODO: record actual environment \\
\bottomrule
\end{tabularx}
\end{center}

These are initial experimental settings, not recovered settings from an existing run. Any revision will be made using training/validation evidence and saved before test evaluation. Hidden-size and dropout ablations will be considered only if the available training sample size supports them. If the experiment is restricted to the small historical query set, learned-gate results will be labeled exploratory and evaluated with grouped out-of-fold predictions rather than in-sample fitting.

\textbf{Reproducibility.} Each run will preserve the split manifest, candidate IDs, relevance mappings, generated-caption provenance, branch rankings, selected weights, gate checkpoint, per-query metrics, timing logs, and resolved configuration. Data preprocessing and caption caches are shared across methods. Reported improvements will be attributed to fusion only when the underlying branch outputs are identical.

